% !TEX root = task-2-11-4.tex
% Editable task bank. Append new BankTask blocks at the end of the document.
\documentclass[12pt,a4paper]{article}
\usepackage{physics-textbook}
\usepackage{graphicx,needspace}
\usetikzlibrary{decorations.pathmorphing}
\newif\ifShowSolutions
\ifdefined\TutorEdition
  \ShowSolutionstrue
\else
  \ShowSolutionsfalse % Change to \ShowSolutionstrue to display solutions.
\fi
\newcounter{banktask}
\newsavebox{\BankTaskBox}
\newcommand{\BankTaskContent}[3]{%
  \noindent\textbf{\thebanktask.}\quad #2\par
  \ifShowSolutions
    \Needspace{50mm}% Keep the solution heading with its opening text or diagram.
    #3\par
    \ifdefined\TutorEdition
      \ifcsname TutorSource#1\endcsname
        \nopagebreak\smallskip{\small\csname TutorSource#1\endcsname\par}%
      \fi
    \fi
  \fi
}
\newcommand{\BankTask}[3]{%
  \refstepcounter{banktask}\par\medskip
  \sbox{\BankTaskBox}{\begin{minipage}{\linewidth}%
    \BankTaskContent{#1}{#2}{#3}\end{minipage}}%
  \ifdim\dimexpr\ht\BankTaskBox+\dp\BankTaskBox\relax>\dimexpr\textheight-20mm\relax
    \Needspace{35mm}\BankTaskContent{#1}{#2}{#3}%
  \else
    \Needspace{\dimexpr\ht\BankTaskBox+\dp\BankTaskBox\relax}%
    \noindent\usebox{\BankTaskBox}\par
  \fi
  \medskip
}
\begin{document}
\BookSubsection{2.11.4}{Первый закон термодинамики}


% Source sdamgia-29029: https://phys-ege.sdamgia.ru/problem?id=29029
\BankTask{sdamgia-29029}{
В некотором процессе идеальный газ получил количество теплоты $50\,\text{Дж}$. При этом внутренняя энергия газа увеличилась на $20\,\text{Дж}$. Чему равна работа, совершенная газом в этом процессе?

(В Дж.)
}{
\smallskip\noindent\textbf{Решение}\par
По первому закону термодинамики $Q=A + \Delta U,$ следовательно:

$A=Q - \Delta U=50 - 20=30 \text{Дж}.$

Ответ: $30$.
}

% Source sdamgia-35991: https://phys-ege.sdamgia.ru/problem?id=35991
\BankTask{sdamgia-35991}{
Внешние силы совершили над газом работу $400\,\text{Дж}$, а внутренняя энергия газа увеличилась при этом на $700\,\text{Дж}$. Какое количество теплоты газ получил от окружающей среды? (В джоулях.)
}{
\smallskip\noindent\textbf{Решение}\par
По первому закону термодинамики: 

$Q = \Delta U - A'=700 - 400 = 300 \text{Дж}.$

Ответ: $300$.
}

% Source sdamgia-42274: https://phys-ege.sdamgia.ru/problem?id=42274
\BankTask{sdamgia-42274}{
Внешние силы совершили над газом работу $60\,\text{Дж}$, при этом внутренняя энергия газа уменьшилась на $40\,\text{Дж}$. Какое количество теплоты отдал газ в этом процессе? (В джоулях.)
}{
\smallskip\noindent\textbf{Решение}\par
По первому закону термодинамики $Q=A + \Delta U= - A' + \Delta U= - 60 - 40= - 100\text{Дж}.$ Таким образом, газ отдал $100\,\text{Дж}$ теплоты.

Ответ: $100$.
}

% Source sdamgia-38605: https://phys-ege.sdamgia.ru/problem?id=38605
\BankTask{sdamgia-38605}{
В ходе некоторого процесса идеальный одноатомный газ получил количество теплоты, равное $900\,\text{Дж}$. На сколько уменьшилась внутренняя энергия газа, если он совершил работу, равную $1500\,\text{Дж}$. (В джоулях.)
}{
\smallskip\noindent\textbf{Решение}\par
По первому закону термодинамики 

$\Delta U = Q - A = 900 - 1500 = - 600\text{Дж}.$

 Внутренняя энергия идеального газа уменьшилась на $600\,\text{Дж}$.

Ответ: $600$.
}

% Source sdamgia-41541: https://phys-ege.sdamgia.ru/problem?id=41541
\BankTask{sdamgia-41541}{
В цилиндре двигателя в ходе адиабатического процесса внутренняя энергия газа уменьшилась на $200\,\text{Дж}$. Какая работа в результате этого совершена газом? (В джоулях.)
}{
\smallskip\noindent\textbf{Решение}\par
При адиабатном процессе $Q=0,$ поэтому по первому закону термодинамики $A= - \Delta U=200\text{Дж}.$

Ответ: $200$.
}

% Source sdamgia-38760: https://phys-ege.sdamgia.ru/problem?id=38760
\BankTask{sdamgia-38760}{
Газ получил от окружающей среды количество теплоты $600\,\text{Дж}$, при этом внешние силы совершили над ним работу, равную $300\,\text{Дж}$. Масса газа не менялась. На сколько увеличилась внутренняя энергия газа? (В джоулях.)
}{
\smallskip\noindent\textbf{Решение}\par
По первому закону термодинамики 

$\Delta U = Q + A' = 600 + 300 = 900 \text{Дж}.$

Ответ: $900$.
}

% Source sdamgia-3517: https://phys-ege.sdamgia.ru/problem?id=3517
\BankTask{sdamgia-3517}{
Для анализа изотермического, изобарного и изохорного процессов над фиксированным количеством идеального газа используют первое начало термодинамики: $Q=A + \Delta U.$ Передаваемое количество теплоты при:

\TaskMatchingColumns{Изопроцесс}{А)  Изотермическом процессе\par Б)  Изобарном процессе\par В)  Изохорном процессе}{Физическое явление}{$1$.  Идет на увеличение его внутренней энергии\par $2$.  Полностью превращается в работу\par $3$.  Идет на увеличение его внутренней энергии и на работу}

\begin{center}\begin{tabular}{|c|c|c|}\hline
$A$ & Б & В\\\hline
\strut & \strut & \strut\\\hline\end{tabular}\end{center}
}{
\smallskip\noindent\textbf{Решение}\par
Вспомним все, что может потребоваться при решении данной задачи. 

Для идеального газа внутренняя энергия зависит только от температуры: $\Delta U= \frac{3}{2} \nu R\Delta T.$ Следовательно, изменить внутреннюю энергию идеального газа можно, только нагрев или охладив его. Работа газа всегда связана с изменением его объема, если газ расширяется, он совершает положительную работу, если сжимается, то отрицательную. Наконец, идеальный газ подчиняется уравнению состояния Клапейрона — Менделеева: $pV=\nu RT.$

При изотермическом процессе температура остается неизменной, поэтому внутренняя энергия идеального газа не изменяется, а значит, согласно первому началу термодинамики, все переданное тепло полностью превращается в работу (А — $2$). При изобарном процессе, то есть при постоянном давлении, согласно уравнению Клапейрона — Менделеева, меняется как температура, так и объем газа, а значит, в правой части первого начала термодинамики отличны от нуля оба слагаемых: тепло идет и на увеличение внутренней энергии газа, и на работу (Б — $3$). Наконец, при изохорном процессе объем газа фиксирована. Следовательно, все тепло идет на изменение внутренней энергии (В — $1$).

Ответ: $231$.
}

% Source sdamgia-10440: https://phys-ege.sdamgia.ru/problem?id=10440
\BankTask{sdamgia-10440}{
\begin{center}\begin{tikzpicture}[task graph,>=stealth]\draw[->](-.55,0)--(5,0)node[below]{\(p\)};
\draw[->](0,-.55)--(0,3.7)node[left]{\(T\)};
\node[below left]at(0,0){\(0\)};
\draw[thin](0,0)--(1,1)(1,0)--(1,1)(3,0)--(3,3);
\fill(1,0)circle(1.1pt)node[below]{\(p\)};
\fill(3,0)circle(1.1pt)node[below]{\(3p\)};
\draw(3.75,1)--(1,1)--(3,3);
\draw[->](2.95,1)--(2.05,1);
\draw[->](1.5,1.5)--(2.16,2.16);
\fill(3.75,1)circle(1.5pt)node[right]{\(1\)};
\fill(1,1)circle(1.5pt)node[above left]{\(2\)};
\fill(3,3)circle(1.5pt)node[above left]{\(3\)};\end{tikzpicture}\end{center}
Один моль идеального одноатомного газа совершает процесс $1$–$2$–$3$, изображенный на рисунке. В процессе $1$–$2$ температура газа постоянна и равна $300$ К. Какое количество теплоты получил газ на участке $2$–$3$?
}{
\smallskip\noindent\textbf{Решение}\par
В процессе $2$–$3$ происходит изохорное нагревание газа. Работа газа равна нулю и, согласно первому началу термодинамики, переданное газу тепло в таком процессе идет только на изменение его внутренней энергии:

\[Q=\Delta U= \frac{3}{2} \nu R \Delta T.\]

При постоянном объеме для идеального газа выполняется закон Шарля, согласно которому $\frac{p}{T} =const.$ Давление на участке $2$–$3$ увеличилось в $3$ раза, а значит, и температура увеличилась в $3$ раза и стала равной $900$ К.
Найдем количество теплоты, которое получил газ:

\[Q= \frac{3}{2} \nu R \Delta T= \frac{3}{2} \cdot 1 \cdot 8{,}31 \cdot 600\approx 7480\,\text{Дж}.\]

Ответ: $Q\approx 7480\,\text{Дж}.$
}

% Source sdamgia-9326: https://phys-ege.sdamgia.ru/problem?id=9326
\BankTask{sdamgia-9326}{
При сжатии идеального одноатомного газа при постоянном давлении внешние силы совершили работу $2000\,\text{Дж}$. Какое количество теплоты было передано при этом газом окружающим телам? (Ответ дайте в джоулях.)
}{
\smallskip\noindent\textbf{Решение}\par
Изменение внутренней энергии идеального одноатомного газа при постоянном давлении равно $\Delta U= \frac{3}{2} \nu R\Delta T,$ работа газа при изобарном процессе $A=\nu R\Delta T.$ Причем, изменение внутренней энергии при изобарном сжатии отрицательное:

\[\Delta U= - \frac{3}{2} A_{\,\text{внеш}}.\]

Количество теплоты, отданное газом окружающим телам, по первому закону термодинамики равно:

\[Q_{\,\text{отд}}=A_{\,\text{внеш}} - \Delta U=A_{\,\text{внеш}} + \frac{3}{2} A_{\,\text{внеш}}= \frac{5}{2} A_{\,\text{внеш}}= \frac{5}{2} \cdot 2000 = 5000\,\text{Дж}.\]

Ответ: $5000\,\text{Дж}$.
}

% Source sdamgia-25257: https://phys-ege.sdamgia.ru/problem?id=25257
\BankTask{sdamgia-25257}{
В ходе адиабатного процесса над идеальным одноатомным газом была совершена работа $1495{,}8\,\text{Дж}$. В результате температура газа повысилась на $40$ °C. Найдите количество вещества этого газа.
}{
\smallskip\noindent\textbf{Решение}\par
По первому закону термодинамики $\Delta U=Q + A'.$ При адиабатном процессе Q  =  $0$. Изменение внутренней энергии для одноатомного газа $\Delta U= \frac{3}{2} v R\Delta T.$ Тогда количество вещества: 

\[v = \frac{2A'}{3R\Delta T} = \frac{2 \cdot 1495{,}8}{3 \cdot 8{,}31 \cdot 40} =3\,\text{моль}.\]

Ответ: $3$ моль.
}


% Source BA6eFF: https://ege.fipi.ru/bank/questions.php?search=1&proj=BA1F39653304A5B041B656915DC36B38&qid=BA6eFF
\BankTask{BA6eFF}{
Двум моль газа сообщили количество теплоты, равное \(1700\ \text{Дж}\), при этом внешние силы совершили над газом работу, равную \(500\ \text{Дж}\). На сколько увеличилась внутренняя энергия газа (в $\text{Дж}$)?
}{
\smallskip\noindent\textbf{Решение от GPT}\par
По первому закону термодинамики
\[
\Delta U=Q+A_{\text{внеш}}=1700+500=2200\ \text{Дж}.
\]
Ответ: \(2200\ \text{Дж}\).
}

% Source 14FBF5: https://ege.fipi.ru/bank/questions.php?search=1&proj=BA1F39653304A5B041B656915DC36B38&qid=14FBF5
\BankTask{14FBF5}{
Газ получил количество теплоты, равное \(300\ \text{Дж}\), при этом внутренняя энергия газа уменьшилась на \(100\ \text{Дж}\). Масса газа не менялась. Какую работу совершил газ в этом процессе (в $\text{Дж}$)?
}{
\smallskip\noindent\textbf{Решение от GPT}\par
По первому закону термодинамики \(Q=\Delta U+A\), где \(A\) — работа газа. Так как \(\Delta U=-100\ \text{Дж}\),
\[
A=Q-\Delta U=300-(-100)=400\ \text{Дж}.
\]
Ответ: \(400\ \text{Дж}\).
}

% Source 5AD1F4: https://ege.fipi.ru/bank/questions.php?search=1&proj=BA1F39653304A5B041B656915DC36B38&qid=5AD1F4
\BankTask{5AD1F4}{
Газ в сосуде сжали, совершив работу, равную \(500\,\text{Дж}\). Внутренняя энергия газа при этом увеличилась на \(350\,\text{Дж}\). Какое количество теплоты отдал газ окружающей среде (в $\text{Дж}$)?
}{
\smallskip\noindent\textbf{Решение от GPT}\par
По первому закону термодинамики \(\Delta U=Q+A_{\text{внеш}}\), где \(Q\) — теплота, полученная газом. Тогда \(Q=350-500=-150\,\text{Дж}\). Газ отдал окружающей среде \(Q_{\text{отд}}=150\,\text{Дж}\).

Ответ: \(150\,\text{Дж}\).
}

% Source 7166: https://phys-ege.sdamgia.ru/problem?id=7166
\BankTask{7166}{
\begin{center}\begin{tikzpicture}[task graph,>=stealth]\draw[->](0,0)--(5,0)node[below]{\(T\)};\draw[->](0,0)--(0,3.7)node[left]{\(p\)};\node[below left]at(0,0){\(0\)};\fill(1,1.15)circle(1.25pt)node[left]{\(1\)};\fill(1.85,2.12)circle(1.25pt)node[above]{\(2\)};\fill(3.6,2.12)circle(1.25pt)node[above]{\(3\)};\draw(1,1.15)--(1.85,2.12);\draw[->](1,1.15)--(1.527,1.751);\draw(1.85,2.12)--(3.6,2.12);\draw[->](1.85,2.12)--(2.935,2.120);\draw[thin](0,0)--(1,1.15);\end{tikzpicture}\end{center}
Один моль одноатомного идеального газа совершает процесс $1$–$2$–$3$, график которого показан на рисунке в координатах p–T. Известно, что давление газа p в процессе $1$–$2$ увеличилось в $2$ раза. Какое количество теплоты было сообщено газу в процессе $1$–$2$–$3$, если его температура T в состоянии $1$ равна $300$ К, а в состоянии $3$ равна $900$ К?
}{
\smallskip\noindent\textbf{Решение}\par
Для определения количества теплоты $Q_{123}$ необходимо сложить количества теплоты, сообщенные газу на участках $1$–$2$ и $2$–$3$: $Q_{123}$  =  $Q_{12}$ + $Q_{23}$.
Исходя из приведенного графика, можно сделать вывод, что процесс $1$–$2$ является изохорным. Для него, как следует из уравнения Клапейрона  — Менделеева, $\frac{p_1}{T_1} = \frac{p_2}{T_2} ,$ откуда $\frac{p_2}{p_1} = \frac{T_2}{T_1} =2.$ Следовательно:

\[T_2=T_1 \frac{p_2}{p_1} =2T_1=300 \cdot 2=600\,\text{К}.\]

Работа газа в процессе $1$–$2$ равна нулю, и для него первый закон термодинамики с учетом выражения для внутренней энергии одноатомного идеального газа принимает вид:

\[Q_{12}=\Delta U_{12}= \frac{3}{2} \nu R ( T_2 - T_1 ) = \frac{3}{2} \nu RT_1.\]

Процесс $2$–$3$ является изобарным с давлением p  =  $p_{2}$  =  const, для него первый закон термодинамики принимает вид: $Q=A_{23} + \Delta U_{23},$ где $\Delta U_{23}= \frac{3}{2} \nu R ( T_3 - T_2 )$  — изменение внутренней энергии газа, $A_{23}=p_2 ( V_3 - V_2 )$  — совершенная газом работа. Из уравнения Клапейрона  — Менделеева $pV=\nu RT$ следует, что:

\[Q_{23}= \frac{3}{2} \nu R ( T_3 - T_2 ) + \nu R ( T_3 - T_2 ) = \frac{5}{2} \nu R ( T_3 - T_2 ) .\]

Таким образом:

\[Q_{23}= \frac{5}{2} \nu R ( T_3 - 2T_1 ) .\]

В результате:

\[Q_{123}= \frac{3}{2} \nu RT_1 + \frac{5}{2} \nu R ( T_3 - 2T_1 ) = \frac{3}{2} \cdot 1 \cdot 8{,}31 \cdot 300 + \frac{5}{2} \cdot 1 \cdot 8{,}31 \cdot ( 900 - 600 ) \approx10\,\text{кДж}.\]

Ответ: $10\,\text{кДж}$.
}

% Source 33046: https://phys-ege.sdamgia.ru/problem?id=33046
\BankTask{33046}{
В вертикальном цилиндре с гладкими стенками под массивным металлическим поршнем находится одноатомный идеальный газ. В начальном состоянии поршень массой M и площадью основания S покоится на высоте h, опираясь на выступы (см. рис. $1$). Давление газа $p_{0}$ равно внешнему атмосферному. Какое количество теплоты Q нужно сообщить газу при медленном его нагревании, чтобы поршень оказался на высоте H (cм. рис. $2$)? Тепловыми потерями пренебречь.

\begin{center}\begin{tikzpicture}[task graph,>=stealth]\draw (0,3.5)--(0,0)--(2.25,0)--(2.25,3.5);\draw(-.1,0)--(2.4,0);\foreach \x in {0,.16,...,2.4}{\draw(\x,0)--++(-.15,-.16);}\fill(0,1.2)--(.23,1.2)--(0,.97)--cycle;\fill(2.25,1.2)--(2.02,1.2)--(2.25,.97)--cycle;\draw[fill=orange!30](0,1.2)rectangle++(2.25,.32);\draw(2.25,1.2)--(3.1,1.2);\draw[<->](2.85,0)--(2.85,1.2)node[midway,right]{\(h\)};\node at(1.05,2.45){\(p_0\)};\node at(1.1,.53){\(p_0\)};\draw(1.8,1.52)--(2.65,2.2)node[right]{\(M,S\)};\draw[->](3.45,1.53)--(4.3,1.53);\end{tikzpicture}\end{center}
Рис. $1$

\begin{center}\begin{tikzpicture}[task graph,>=stealth]\draw (0,3.5)--(0,0)--(2.25,0)--(2.25,3.5);\draw(-.1,0)--(2.4,0);\foreach \x in {0,.16,...,2.4}{\draw(\x,0)--++(-.15,-.16);}\fill(0,1.2)--(.23,1.2)--(0,.97)--cycle;\fill(2.25,1.2)--(2.02,1.2)--(2.25,.97)--cycle;\draw[fill=orange!30](0,2.35)rectangle++(2.25,.32);\draw(2.25,2.35)--(3.1,2.35);\draw[<->](2.85,0)--(2.85,2.35)node[midway,right]{\(H\)};\node at(1.2,3.15){\(p_0\)};\end{tikzpicture}\end{center}
Рис. $2$
}{
\smallskip\noindent\textbf{Решение}\par
Изобразим силы, действующие на поршень в процессе движения.
\begin{center}\begin{tikzpicture}[task graph,>=stealth]\draw[fill=orange!30](-.78,-.25)rectangle(.78,.25);\fill (0,0)circle(1.3pt);\draw[->](0,0)--(0,1.55)node[right]{\(p_2S\)};\draw[->](0,0)--(0,-.77)node[right]{\(Mg\)};\draw[->](0,0)--(0,-1.4)node[right]{\(p_0S\)};\end{tikzpicture}\end{center}
До момента, пока $p_2S < p_0S + Mg,$ поршень будет покоиться, давление газа увеличиваться. При достижении равновесия $p_2S=p_0S + Mg$ поршень начнет двигаться вверх, так как нагревание медленное, из-за чего процесс можно считать изобарным.
По первому закону термодинамики $Q=A + \Delta U.$
Изменение внутренней энергии равно $\Delta U= \frac{3}{2} \nu R\Delta T.$
Работа газа при изобарном расширении:

\[A=p_2\Delta V= ( p_0 + \frac{Mg}{S} ) S ( H - h ) = ( p_0S + Mg ) ( H - h ) .\]

Пусть начальный объем $V_0=Sh,$ конечный объем $V=SH.$
Запишем уравнение Клапейрона-Менделеева для каждого состояния газа:

$p_0V_0=\nu RT_0;$
$p_2V_=\nu RT_2.$

При объединении уравнений, получаем:

\[( p_0S + Mg ) H - p_0Sh=\nu R\Delta T.\]

Тогда изменение внутренней энергии газа равно:

\[\Delta U= \frac{3}{2} ( ( p_0S + Mg ) H - p_0Sh ) .\]

Тогда полученное газом количество теплоты равно:

$Q= \frac{3}{2} ( ( p_0S + Mg ) H - p_0Sh ) + ( p_0S + Mg ) ( H - h ) ;$
$Q= 2{,}5p_0S ( H - h ) + Mg ( 2{,}5H - h ) .$

Ответ: $Q= 2{,}5p_0S ( H - h ) + Mg ( 2{,}5H - h ) .$
}

% Source 7875: https://phys-ege.sdamgia.ru/problem?id=7875
\BankTask{7875}{
\begin{center}\begin{tikzpicture}[task graph,>=stealth]\draw(5,0)--(0,0)--(0,2.65)--(5,2.65);\draw[fill=orange!25](2.26,0)rectangle(2.66,2.65);\foreach \y in {.13,.3,...,2.63}{\draw(2.26,\y)--++(.4,-.12);}\draw(0,1.26)--(.18,1.26);\foreach \x in {.2,.45,...,1.95}{\draw (\x,1.26)arc(180:-145:.13 and .19);}\draw(2.05,1.26)--(2.26,1.26);\node at(1.05,1.91){\(k\)};\node at(.96,.55){\(p_0\)};\node at(4,1.15){\(p_0\)};\draw(2.63,2.1)--(3.8,3.14)node[right]{\(S\)};\draw(0,0)--(0,-.95)(2.26,0)--(2.26,-.95);\draw( .22,-.74)--(2.04,-.74)node[midway,above]{\(L\)};\end{tikzpicture}\end{center}
В горизонтальном цилиндре с гладкими стенками под массивным поршнем с площадью S находится одноатомный идеальный газ. Поршень соединен с основанием цилиндра пружиной с жесткостью k. В начальном состоянии расстояние между поршнем и основанием цилиндра равно L, а давление газа в цилиндре равно внешнему атмосферному давлению $p_{0}$ (см. рис.). Какое количество теплоты Q передано затем газу, если в результате поршень медленно переместился вправо на расстояние b?
}{
\smallskip\noindent\textbf{Решение}\par
Тепло, переданное газу, идет на изменение его внутренней энергии и на совершением им работы:

\[Q=\Delta U + A.\]

В начальном состоянии давление и объем газа равны $p_0$ и $V_0=SL.$ Следовательно, пружина была не деформирована. 
В конечном состоянии пружина была деформирована на величину b, возникла сила упругости $F_{\,\text{упр}}=kb,$ направленная влево. Запишем условие равновесия поршня в конечном состоянии:

\[p_0S + F_{\,\text{упр}}=p_{\,\text{к}}S.\]

Тогда давление воздуха в конечном состоянии $p_{\,\text{к}}=p_0 + \frac{kb}{S}$ и объем воздуха в конечном состоянии  $V_{\,\text{к}}=S ( L + b ) .$ 
Изменение внутренней энергии воздуха $\Delta U = \frac{3}{2} \nu R\Delta T.$ Тогда получаем:

\[\Delta U = \frac{3}{2} ( p_{\,\text{к}}V_{\,\text{к}} - p_0V_0 ) = \frac{3}{2} ( p_0Sb + kbL + kb ^2 ) .\]

Чтобы рассчитать работу, заметим, что в каждый момент времени, когда поршень сдвинут на x от начального положения давление равно $p=p_0 + \frac{kx}{S} =p_0 + \frac{k ( V - V_0 )}{S ^2} ,$ то есть давление линейно зависит от объема. Значит, на pV-диаграмме процесс расширения будет изображен отрезком прямой, а фигура под графиком будет являться трапецией, площадь которой равна:

\[A = \frac{p_0 + p_{\,\text{к}}}{2} ( V_{\,\text{к}} - V_0 ) =p_0Sb + \frac{kb ^2}{2} .\]

Таким образом, полученное воздухом количество теплоты равно:

\[Q = \frac{5}{2} p_0Sb + \frac{3}{2} kbL + 2kb ^2 .\]

Ответ: $Q = \frac{5}{2} p_0Sb + \frac{3}{2} kbL + 2kb ^2 .$
}

% Source 7129: https://phys-ege.sdamgia.ru/problem?id=7129
\BankTask{7129}{
\begin{center}\begin{tikzpicture}[task graph,>=stealth]\draw(4.7,0)--(0,0)--(0,2.65)--(4.7,2.65);\draw[fill=orange!35](2.5,0)rectangle(2.94,2.65);\draw[dashed](3.95,0)rectangle(4.39,2.65);\node at(1.3,1.62){\(p\)};\draw(2.82,1.02)--(3.57,1.94)node[above]{\(S\)};\draw(0,0)--(0,-.65)(2.5,0)--(2.5,-.91)(3.95,0)--(3.95,-.91);\draw[<->](0,-.49)--(2.5,-.49)node[midway,below]{\(L\)};\draw[<->](2.5,-.82)--(3.95,-.82)node[midway,below]{\(x\)};\end{tikzpicture}\end{center}
В горизонтальном цилиндрическом сосуде, закрытом поршнем, находится одноатомный идеальный газ. Первоначальное давление газа p  =  $4$ · $10^{5}$ Па. Расстояние от дна сосуда до поршня равно L. Площадь поперечного сечения поршня S  =  $25\,\text{см}^{2}$. В результате медленного нагревания газ получил количество теплоты Q  =  $1{,}65\,\text{кДж}$, а поршень сдвинулся на расстояние x  =  $10\,\text{см}$. При движении поршня на него со стороны стенок сосуда действует сила трения величиной $F_{\,\text{тр}}$  =  $3$ · $10^{3}$ Н. Найдите L. Считать, что сосуд находится в вакууме.
}{
\smallskip\noindent\textbf{Решение}\par
$1$.  Поршень будет медленно двигаться, если сила давления газа на поршень и сила трения со стороны стенок сосуда уравновесят друг друга: $p_{2}$S = $F_{\,\text{тр}}$, откуда $p_2= \frac{F_{\,\text{тр}}}{S} = 12 \cdot 10 ^{5} \,\text{Па} > p_1.$ $( p_1 = p ) .$
$2$.  Поэтому при нагревании газа поршень будет неподвижен, пока давление газа не достигнет значения $p_{2}$. В этом процессе газ получает количество теплоты $Q_{12}$. Затем поршень будет сдвигаться, увеличивая объем газа, при постоянном давлении. В этом процессе газ получает количество теплоты $Q_{23}$.
$3$.  В процессе нагревания, в соответствии с первым началом термодинамики, газ получит количество теплоты:

\[Q=Q_{12} + Q_{23} = ( U_3 - U_1 ) + p_2Sx = ( U_3 - U_1 ) + F_{\,\text{тр}}x.\]

$4$.  Внутренняя энергия одноатомного идеального газа: 

\[U_1= \frac{3}{2} \nu RT_1= \frac{3}{2} p_1SL\]

в начальном состоянии, 

\[U_3= \frac{3}{2} \nu RT_3= \frac{3}{2} p_2S ( L + x ) = \frac{3}{2} F_{\,\text{тр}} ( L + x )\]

в конечном состоянии.
$5$.  Из пп. $3$, $4$ получаем $L= \frac{Q - \frac{5}{2} F_{\,\text{тр}}x}{\frac{3}{2} ( F_{\,\text{тр}} - p_1S )} = \frac{1650 - 2{,}5 \cdot 3000 \cdot 0{,}1}{1{,}5 \cdot ( 3000 - 400000 \cdot 0{,}0025 )} =0{,}3\,\text{м}.$

Ответ: L  =  $0{,}3\,\text{м}$.
}

\end{document}
